% Bien dich bang: xelatex quy_trinh_chuan_bi_du_lieu.tex (chay 2 lan)
\documentclass[12pt,a4paper]{article}

\usepackage{fontspec}
\usepackage[vietnamese]{babel}
\defaultfontfeatures{Ligatures=TeX}
\setmainfont{TeX Gyre Termes}
\setsansfont{TeX Gyre Heros}
\setmonofont{DejaVu Sans Mono}[Scale=0.85]

\usepackage[a4paper,margin=2.5cm]{geometry}
\usepackage{microtype}
\usepackage{setspace}
\usepackage{booktabs}
\usepackage{longtable}
\usepackage{array}
\usepackage{enumitem}
\usepackage{xcolor}
\usepackage{hyperref}
\usepackage{tikz}
\usetikzlibrary{arrows.meta,positioning,shapes.geometric}

\hypersetup{
  colorlinks=true,
  linkcolor=blue!50!black,
  urlcolor=blue!60!black,
  pdftitle={Quy trình chuẩn bị dữ liệu và xử lý nhãn cấu trúc},
  pdfauthor={Nhóm 4}
}
\setstretch{1.15}
\setlength{\parindent}{1.25cm}
\setlength{\parskip}{0.25em}
\renewcommand{\arraystretch}{1.15}

\newcommand{\code}[1]{\texttt{#1}}
\newcommand{\nullvalue}{\texttt{null}}

\title{\textbf{QUY TRÌNH CHUẨN BỊ DỮ LIỆU\\VÀ XỬ LÝ NHÃN CẤU TRÚC}}
\author{Nhóm 4}
\date{Tháng 8 năm 2026}

\begin{document}
\maketitle

\begin{abstract}
Báo cáo trình bày quy trình chuẩn bị tập dữ liệu gồm tám văn bản quy phạm pháp
luật trong lĩnh vực lao động và bảo hiểm xã hội. Quy trình bao gồm lựa chọn nguồn
dữ liệu, trích xuất lớp văn bản hoặc nhận dạng ký tự quang học, làm sạch nội dung,
khôi phục cấu trúc pháp lý, xử lý nhãn cấu trúc bị nhiễu, chia đoạn theo cấu trúc
và gắn siêu dữ liệu. Kết quả thu được 668 đoạn dữ liệu ở định dạng JSON Lines từ
476 trang PDF. Toàn bộ tám văn bản đều có dữ liệu, không có định danh đoạn trùng
lặp và không có đoạn rỗng.
\end{abstract}

\tableofcontents
\newpage

\section{Mục tiêu và phạm vi}

Mục tiêu của bước chuẩn bị dữ liệu là chuyển các văn bản pháp luật ở định dạng
PDF thành tập dữ liệu có cấu trúc, phù hợp cho các bước phân tích hoặc huấn luyện
mô hình về sau. Khác với tài liệu thông thường, ý nghĩa của một đoạn văn bản luật
phụ thuộc mạnh vào vị trí của nó trong Phần, Chương, Mục, Điều, Khoản và Điểm.
Vì vậy, quy trình không chia văn bản hoàn toàn theo độ dài cố định mà ưu tiên bảo
toàn đơn vị pháp lý và ngữ cảnh phân cấp.

Tập đầu vào gồm đúng tám tệp PDF được đánh số từ 01 đến 08 trong thư mục
\code{data}. Ba tài liệu trong thư mục \code{data/tuychon} không được sử dụng.
Danh sách đầu vào được khai báo tường minh thay vì quét thư mục đệ quy, nhờ đó
kết quả có thể tái lập và không bị thay đổi khi xuất hiện tệp ngoài phạm vi.

Trong báo cáo này, ``nhãn'' được hiểu là nhãn cấu trúc và nhãn siêu dữ liệu gắn
cho đơn vị nội dung, chẳng hạn số Điều, danh sách Khoản, danh sách Điểm, loại đoạn
và vị trí trang nguồn. Đây không phải là nhãn dự đoán nội dung pháp lý do con
người gán thủ công.

\section{Tổng quan quy trình}

Quy trình được tổ chức thành các bước tuần tự như Hình~\ref{fig:pipeline}.
Thông tin trang nguồn được duy trì xuyên suốt các bước để có thể đối chiếu một
đoạn JSONL với PDF ban đầu.

\begin{figure}[htbp]
  \centering
  \begin{tikzpicture}[
    node distance=7mm,
    box/.style={rectangle,rounded corners,draw=blue!60!black,fill=blue!5,
      text width=10.5cm,align=center,minimum height=9mm},
    arrow/.style={-{Latex[length=2.5mm]},thick,draw=blue!60!black}
  ]
    \node[box] (input) {Danh sách cố định 8 tệp PDF};
    \node[box,below=of input] (extract) {Trích xuất text layer hoặc OCR tiếng Việt theo trang};
    \node[box,below=of extract] (clean) {Loại nhiễu và chuẩn hóa Unicode, khoảng trắng, dấu câu};
    \node[box,below=of clean] (parse) {Khôi phục Phần--Chương--Mục--Điều--Khoản--Điểm};
    \node[box,below=of parse] (label) {Chuẩn hóa và kiểm tra nhãn cấu trúc};
    \node[box,below=of label] (chunk) {Structure-aware chunking và gắn metadata};
    \node[box,below=of chunk] (output) {JSONL, dữ liệu trung gian và báo cáo QA};
    \draw[arrow] (input) -- (extract);
    \draw[arrow] (extract) -- (clean);
    \draw[arrow] (clean) -- (parse);
    \draw[arrow] (parse) -- (label);
    \draw[arrow] (label) -- (chunk);
    \draw[arrow] (chunk) -- (output);
  \end{tikzpicture}
  \caption{Luồng chuẩn bị dữ liệu và xử lý nhãn cấu trúc.}
  \label{fig:pipeline}
\end{figure}

\section{Chuẩn bị và trích xuất dữ liệu}

\subsection{Kiểm tra lớp văn bản}

Mỗi PDF được kiểm tra số trang và lượng ký tự có thể trích xuất. Kết quả cho thấy
hai tài liệu 03 và 05 có lớp văn bản, trong khi sáu tài liệu 01, 02, 04, 06, 07
và 08 là bản quét thuần ảnh. Do đó, việc chỉ sử dụng công cụ đọc PDF sẽ làm mất
hoàn toàn nội dung của phần lớn tập dữ liệu.

Với hai PDF có lớp văn bản, nội dung được trích xuất trực tiếp bằng \code{pypdf}.
Với sáu PDF ảnh, từng ảnh trang được đưa vào Tesseract sử dụng mô hình tiếng Việt
\code{vie.traineddata}. OCR chỉ đóng vai trò fallback trên trang không có chữ,
nhờ đó không làm suy giảm chất lượng của các trang đã có lớp văn bản tốt. Kết quả
được lưu theo từng trang với ba thành phần: số trang PDF, văn bản thô và danh sách
dòng theo đúng thứ tự xuất hiện.

\subsection{Kiểm soát lỗi trích xuất}

Một tài liệu bị xem là lỗi nếu cả trích xuất trực tiếp và OCR không tạo được nội
dung đủ dùng. Pipeline dừng thay vì âm thầm sinh tệp rỗng. Với mỗi trang, số trang
PDF được giữ nguyên để các bước làm sạch, phân tích cấu trúc và chia đoạn có thể
truyền tiếp phạm vi \code{page\_start}--\code{page\_end}.

\section{Làm sạch và chuẩn hóa văn bản}

\subsection{Loại bỏ nhiễu trình bày}

Header và footer được phát hiện từ các dòng ngắn xuất hiện lặp lại tại ba dòng
không rỗng đầu hoặc cuối của nhiều trang. Một dòng chỉ được xem là header/footer
nếu xuất hiện ở ít nhất 60\% số trang của tài liệu. Các số trang đứng độc lập ở
vùng lề cũng được loại bỏ. Mỗi dòng bị xóa được ghi lại cùng số trang, vị trí,
nội dung gốc và nguyên nhân, giúp quá trình làm sạch có thể kiểm tra lại.

Quy tắc loại nhiễu được thiết kế theo hướng bảo thủ. Số hiệu văn bản, ngày tháng,
tỷ lệ phần trăm, số tiền, số Khoản và số Điểm không bị loại chỉ vì chúng chứa
chữ số hoặc ký hiệu. Nội dung không khớp mẫu cấu trúc cũng không bị xóa tự động.

\subsection{Chuẩn hóa ký tự và khoảng trắng}

Văn bản được chuẩn hóa Unicode NFC, thay khoảng trắng không ngắt bằng khoảng
trắng thông thường, loại ký tự điều khiển và gộp khoảng trắng liên tiếp. Khoảng
trắng không cần thiết trước dấu câu được loại bỏ, nhưng dấu phẩy thập phân như
\code{5,5\%} vẫn được giữ nguyên. Các dòng được ghép ở bước tạo đoạn, đồng thời
ranh giới trang vẫn được lưu bằng metadata.

\section{Khôi phục cấu trúc và xử lý nhãn}

\subsection{Nhận diện cấu trúc phân cấp}

Bộ phân tích đọc các dòng theo thứ tự và duy trì trạng thái cấu trúc hiện hành.
Khi gặp nhãn \emph{Phần}, \emph{Chương}, \emph{Mục} hoặc \emph{Tiểu mục}, trạng
thái tương ứng được cập nhật. Khi một \emph{Điều} mới xuất hiện, Điều trước đó
được kết thúc và nhận bản sao của đường dẫn cấu trúc hiện hành. Các dòng bắt đầu
bằng số thứ tự được nhận diện là Khoản; các dòng bắt đầu bằng chữ cái và dấu
``)'' hoặc ``.'' được nhận diện là Điểm.

Nội dung trước Điều đầu tiên được gắn nhãn \code{preamble}. Nội dung không nhận
diện được cấp cấu trúc được giữ trong đơn vị hợp lệ gần nhất thay vì bị bỏ đi.
Nhờ vậy, quá trình tạo nhãn không làm giảm độ bao phủ nội dung.

\subsection{Các nhãn cấu trúc được thiết lập}

Các nhãn cấu trúc được tạo theo hai mức: đường dẫn phân cấp của văn bản và đơn
vị nội dung bên trong Điều. Bảng~\ref{tab:structure-labels} trình bày đúng các
nhãn được cài đặt trong parser của dự án.

\begin{longtable}{>{\ttfamily}p{2.8cm}p{3.0cm}p{8.0cm}}
\caption{Các nhãn cấu trúc trong dự án.}\label{tab:structure-labels}\\
\toprule
\normalfont\textbf{Nhãn} & \textbf{Kiểu/giá trị} & \textbf{Ý nghĩa và cách thiết lập} \\
\midrule
\endfirsthead
\toprule
\normalfont\textbf{Nhãn} & \textbf{Kiểu/giá trị} & \textbf{Ý nghĩa và cách thiết lập} \\
\midrule
\endhead
part & object hoặc null & Phần hiện hành, gồm \code{number} và \code{title}. Khi mở Phần mới, nhãn Chương và Mục cũ được đặt lại. \\
chapter & object hoặc null & Chương hiện hành, gồm \code{number} và \code{title}. Số Chương chấp nhận số La Mã hoặc số Ả Rập; khi mở Chương mới, nhãn Mục cũ được đặt lại. \\
section & object hoặc null & Mục hoặc Tiểu mục hiện hành, gồm \code{number} và \code{title}. Hai cấp này được ánh xạ vào cùng trường \code{section}. \\
article & object hoặc null & Điều gốc, gồm \code{number} và \code{title}. Mọi chunk con của Điều đều giữ cùng object này; nội dung ngoài Điều nhận \nullvalue{}. \\
clause\_ids & list[string] & Danh sách số Khoản thực sự xuất hiện trong chunk, nhận diện từ mẫu số Ả Rập ở đầu dòng như \code{1.} hoặc \code{1)}. \\
point\_ids & list[string] & Danh sách ký hiệu Điểm thực sự xuất hiện trong chunk, nhận diện từ chữ cái đầu dòng như \code{a)} hoặc \code{a.}; ký hiệu được chuyển về chữ thường. \\
unit\_type & article, preamble & Nhãn nội bộ trước bước chunking. \code{article} là một Điều; \code{preamble} là nội dung có ý nghĩa trước Điều đầu tiên. \\
\bottomrule
\end{longtable}

Các object \code{part}, \code{chapter}, \code{section} và \code{article} sử dụng
cùng cấu trúc \code{\{"number": ..., "title": ...\}}. Khi một cấp không tồn
tại trong văn bản, trường tương ứng được đặt là \nullvalue{}, không suy diễn một
nhãn giả. \code{clause\_ids} và \code{point\_ids} là danh sách rỗng nếu chunk
không chứa Khoản hoặc Điểm được nhận diện.

\subsection{Chuẩn hóa nhãn bị nhiễu do OCR}

OCR có thể làm sai nhãn cấu trúc, ví dụ \code{Điều} thành \code{Điền}, hoặc đưa
ký tự lề vào trước nhãn như \code{| Điều 14}, \code{\_ Điều 28} và
\code{l Tủ Điều 23}. Bộ phân tích áp dụng nhận diện dung sai cho các biến thể đã
quan sát và cho phép tối đa tám ký tự nhiễu trước từ khóa Điều. Metadata vẫn được
chuẩn hóa về số Điều đúng, còn văn bản OCR trong trường \code{text} được giữ lại
để không che giấu chất lượng nguồn. Mọi trường hợp chuẩn hóa nhãn đều được ghi
thành cảnh báo trong báo cáo QA.

\subsection{Loại nhãn Điều giả}

Một câu viện dẫn có thể bắt đầu bằng cụm \code{Điều 151 và khoản 4 Điều 177...}
nhưng không phải tiêu đề Điều. Để hạn chế gán nhãn sai, một ứng viên tiêu đề phải
có dấu phân cách ``.'', ``:'' hoặc ``-'' ngay sau số Điều. Sau ứng viên đầu tiên,
số Điều phải tăng theo trình tự hợp lý; ứng viên lùi số hoặc nhảy xa bị giữ như
nội dung của Điều hiện hành và được ghi cảnh báo \code{Rejected Article-like
reference}. Cách xử lý này đồng thời loại được các mẫu hợp đồng hoặc quyết định
đính kèm có hệ thống Điều riêng, tránh trộn nhãn của phụ lục vào cấu trúc chính.

\section{Chia đoạn theo cấu trúc và gắn metadata}

Đơn vị chia ưu tiên là một Điều hoàn chỉnh. Ngưỡng mặc định là 6.000 ký tự.
Nếu một Điều vượt ngưỡng, các Khoản liên tiếp được đóng gói vào các đoạn riêng.
Nếu một Khoản vẫn quá dài, ranh giới Điểm được ưu tiên; ranh giới câu chỉ được
dùng khi đơn vị nhỏ hơn vẫn vượt ngưỡng. Tiêu đề Điều và metadata Điều gốc được
giữ trong mọi đoạn con.

\subsection{Schema metadata của mỗi chunk}

Mỗi dòng JSONL là một object độc lập và luôn phát ra đầy đủ các trường trong
Bảng~\ref{tab:metadata-schema}. Trường không có giá trị dùng \nullvalue{} hoặc
danh sách rỗng thay vì bị lược bỏ, nhờ đó schema nhất quán giữa tám văn bản.

\begin{longtable}{>{\ttfamily}p{3.5cm}p{2.7cm}p{7.6cm}}
\caption{Các trường metadata trong một dòng JSONL.}\label{tab:metadata-schema}\\
\toprule
\normalfont\textbf{Trường} & \textbf{Kiểu dữ liệu} & \textbf{Mô tả} \\
\midrule
\endfirsthead
\toprule
\normalfont\textbf{Trường} & \textbf{Kiểu dữ liệu} & \textbf{Mô tả} \\
\midrule
\endhead
chunk\_id & string & Định danh ổn định và duy nhất của chunk, lấy 20 ký tự hexa đầu của SHA-256. \\
document\_id & string & Định danh tài liệu từ \code{doc-01} đến \code{doc-08}. \\
title & string & Tên đầy đủ của văn bản pháp luật. \\
document\_number & string & Số, ký hiệu chính thức của văn bản, ví dụ \code{45/2019/QH14}. \\
document\_type & string & Loại văn bản: \code{bộ luật}, \code{luật}, \code{nghị định} hoặc \code{thông tư}. \\
source\_file & string & Đường dẫn tương đối đến PDF nguồn; dùng để truy vết và kiểm tra phạm vi dữ liệu. \\
page\_start & integer & Trang PDF đầu tiên chứa nội dung chunk, đánh số từ 1. \\
page\_end & integer & Trang PDF cuối cùng chứa nội dung chunk; luôn lớn hơn hoặc bằng \code{page\_start}. \\
part & object/null & Nhãn Phần gồm \code{number}, \code{title}; \nullvalue{} nếu không có. \\
chapter & object/null & Nhãn Chương gồm \code{number}, \code{title}; \nullvalue{} nếu không có. \\
section & object/null & Nhãn Mục/Tiểu mục gồm \code{number}, \code{title}; \nullvalue{} nếu không có. \\
article & object/null & Nhãn Điều gốc gồm \code{number}, \code{title}; \nullvalue{} đối với \code{preamble}. \\
clause\_ids & list[string] & Các số Khoản có mặt trong riêng chunk này, theo thứ tự xuất hiện và không lặp. \\
point\_ids & list[string] & Các ký hiệu Điểm có mặt trong riêng chunk này, theo thứ tự xuất hiện và không lặp. \\
chunk\_type & string & Nhãn loại chunk: \code{preamble}, \code{article} hoặc \code{clause\_group}. \\
chunk\_index & integer & Vị trí chunk trong tài liệu, bắt đầu từ 1 và tăng liên tục. \\
article\_chunk\_index & integer/null & Vị trí chunk con trong Điều gốc; \nullvalue{} nếu nội dung không thuộc Điều. \\
text & string & Nội dung đã làm sạch và chuẩn hóa dùng cho mô hình hoặc truy hồi. \\
char\_count & integer & Số ký tự thực tế của trường \code{text}. \\
token\_estimate & integer & Ước lượng nhanh bằng $\lceil\texttt{char\_count}/3\rceil$; không phải kết quả tokenizer của một mô hình cụ thể. \\
split\_reason & string/null & Lý do chia: \nullvalue{}, \code{max\_size\_clause} hoặc \code{max\_size\_sentence}. \\
processing\_version & string & Phiên bản pipeline đã tạo chunk; bộ dữ liệu hiện tại sử dụng \code{1.0.0}. \\
\bottomrule
\end{longtable}

\subsection{Ý nghĩa các nhãn chunk}

\begin{longtable}{>{\ttfamily}p{3.8cm}p{10.0cm}}
\caption{Tập giá trị nhãn chunk đang được triển khai.}\label{tab:chunk-labels}\\
\toprule
\normalfont\textbf{Giá trị nhãn} & \textbf{Điều kiện gán} \\
\midrule
\endfirsthead
\toprule
\normalfont\textbf{Giá trị nhãn} & \textbf{Điều kiện gán} \\
\midrule
\endhead
preamble & Nội dung trước Điều đầu tiên; \code{article} và \code{article\_chunk\_index} nhận \nullvalue{}. \\
article & Một Điều được giữ nguyên vì không vượt ngưỡng kích thước; \code{split\_reason} nhận \nullvalue{}. \\
clause\_group & Một phần của Điều dài sau khi được chia thành các nhóm; metadata Điều gốc được lặp lại để giữ ngữ cảnh. \\
max\_size\_clause & Giá trị của \code{split\_reason} khi Điều dài được chia theo ranh giới các Khoản. \\
max\_size\_sentence & Giá trị fallback của \code{split\_reason} khi một nhóm đơn lẻ vẫn quá dài và phải chia theo câu. \\
null & Không có thao tác chia; được biểu diễn bằng JSON \nullvalue{}, không phải chuỗi ký tự ``null''. \\
\bottomrule
\end{longtable}

Định danh \code{chunk\_id} được tạo từ SHA-256 của định danh tài liệu, đơn vị
cấu trúc, thứ tự đoạn con và nội dung; 20 ký tự hexa đầu tiên được sử dụng. Vì
vậy, cùng dữ liệu và cấu hình sẽ sinh cùng định danh.

\section{Kết quả và kiểm soát chất lượng}

\begin{longtable}{>{\ttfamily}lrrrrr}
\caption{Thống kê kết quả theo tài liệu.}\label{tab:results}\\
\toprule
ID & Trang & Ký tự sạch & Điều & Chunk & Dòng loại \\
\midrule
\endfirsthead
\toprule
ID & Trang & Ký tự sạch & Điều & Chunk & Dòng loại \\
\midrule
\endhead
doc-01 & 83  & 191,045 & 220 & 226 & 0  \\
doc-02 & 122 & 235,737 & 115 & 132 & 23 \\
doc-03 & 19  & 33,336  & 12  & 19  & 18 \\
doc-04 & 14  & 29,658  & 5   & 10  & 1  \\
doc-05 & 88  & 187,010 & 141 & 148 & 0  \\
doc-06 & 52  & 116,614 & 45  & 49  & 5  \\
doc-07 & 15  & 15,728  & 9   & 11  & 5  \\
doc-08 & 83  & 179,182 & 64  & 73  & 65 \\
\midrule
\textbf{Tổng} & \textbf{476} & \textbf{988,310} & \textbf{611} & \textbf{668} & \textbf{117} \\
\bottomrule
\end{longtable}

Kiểm tra tự động trên đầu ra xác nhận đủ tám định danh tài liệu, không có đường
dẫn thuộc \code{data/tuychon}, không có \code{chunk\_id} trùng và không có đoạn
rỗng. Số Điều trong từng tài liệu tạo thành chuỗi liên tục: 1--220, 1--115,
1--12, 1--5, 1--141, 1--45, 1--9 và 1--64. Chỉ số đoạn con trong mỗi Điều cũng
liên tục từ 1. Bộ kiểm thử gồm 22 trường hợp đã được chạy thành công.

Đầu ra được ghi qua tệp tạm và chỉ thay thế tệp đích sau khi tất cả dòng JSONL
đọc lại hợp lệ. Ngoài tệp tổng hợp \code{legal\_chunks.jsonl}, quy trình còn sinh
tệp riêng cho từng tài liệu, dữ liệu thô/sạch theo trang và báo cáo QA ở định
dạng JSON và Markdown.

\section{Hạn chế của dữ liệu}

Việc chuẩn hóa nhãn cấu trúc không đồng nghĩa với sửa toàn bộ lỗi OCR trong nội
dung. Một số từ vẫn có thể bị nhận sai, chẳng hạn \code{tổ chức} thành
\code{tô chức}, hoặc xuất hiện ký tự lề như \code{¬}. Quy trình chủ động giữ nội
dung OCR và chỉ chuẩn hóa biểu diễn metadata cần thiết cho phân đoạn. Cách làm
này bảo đảm khả năng truy vết nhưng đòi hỏi bước rà soát hoặc sửa lỗi chính tả
riêng nếu dữ liệu được dùng cho tác vụ nhạy cảm với từng ký tự.

Quy tắc tuần tự số Điều phù hợp với tám văn bản chính trong phạm vi đề tài. Khi
áp dụng cho phụ lục có hệ thống Điều độc lập, văn bản sửa đổi có Điều mang hậu
tố phức tạp, hoặc tập trích đoạn không bắt đầu từ cấu trúc chính, cần điều chỉnh
trạng thái parser và metadata loại tài liệu.

\section{Kết luận}

Quy trình đã chuyển đổi tám PDF không đồng nhất về lớp văn bản thành tập JSONL
có cấu trúc và khả năng truy vết. Việc kết hợp trích xuất trực tiếp với OCR theo
trang giúp bao phủ toàn bộ dữ liệu; parser dung sai giúp phục hồi nhãn Điều bị
nhiễu nhưng vẫn loại các câu viện dẫn và cấu trúc phụ dễ gây gán nhãn sai.
Structure-aware chunking bảo toàn ngữ cảnh pháp lý tốt hơn chia đoạn cố định,
đồng thời metadata thống nhất tạo điều kiện cho tìm kiếm, huấn luyện và đánh giá
ở các bước tiếp theo.

\end{document}
